\documentclass[../../main.tex]{subfiles}
\begin{document}
\section{连通性与紧性}\label{section:appA-connectedness-compactness}

如果拓扑空间 $X$ 有两个不相交的非空开子集其并为 $X$，就称它\emph{不连通}；否则称它\emph{连通}。等价地，$X$ 连通，当且仅当 $X$ 中既开又闭的子集只有 $\varnothing$ 和 $X$ 本身。如果 $X$ 是任意拓扑空间，$X$ 的一个\emph{连通子集}是这样的子集：它在赋予子空间拓扑后是连通空间。例如，$\mathbb R$ 的非空连通子集是单点集（一元集合）和区间，即含多于一点、且具有“只要 $a,b\in J$ 且 $a<c<b$ 就有 $c\in J$”这一性质的子集 $J\subset\mathbb R$。

$X$ 的一个极大连通子集（即不被任何更大连通子集真包含的连通子集）称为 $X$ 的一个\emph{分支}（或连通分支）。

\begin{proposition}[连通空间的性质]\label{proposition:appA-connected-properties}
设 $X$ 与 $Y$ 是拓扑空间。
\begin{enumerate}[(a)]
\item 如果 $F:X\to Y$ 连续且 $X$ 连通，则 $F(X)$ 连通。
\item $X$ 的每个连通子集都含于 $X$ 的单个分支中。
\item $X$ 中若干有公共点的连通子空间之并连通。
\item $X$ 的分支是不相交的非空闭子集，其并为 $X$，因而它们构成 $X$ 的一个划分。
\item 如果 $S$ 是 $X$ 中既开又闭的子集，则 $S$ 是 $X$ 若干分支之并。
\item 连通空间的任意有限积连通。
\item 连通空间的任意商空间连通。
\end{enumerate}
\end{proposition}

\begin{exercise}\label{exercise:appA-A-40}
证明前一个命题。
\end{exercise}

与连通性密切相关的是道路连通性。如果 $X$ 是拓扑空间而 $p,q\in X$，$X$ 中从 $p$ 到 $q$ 的一条\emph{道路}是连续映射 $f:I\to X$（其中 $I=[0,1]$）满足 $f(0)=p$、$f(1)=q$。如果对每对点 $p,q\in X$ 都存在 $X$ 中从 $p$ 到 $q$ 的道路，则称 $X$ \emph{道路连通}。$X$ 的\emph{道路分支}是其极大道路连通子集。

\begin{proposition}[道路连通空间的性质]\label{proposition:appA-path-connected-properties}
\begin{enumerate}[(a)]
\item 命题 A.39 中把“连通”处处换成“道路连通”、把“分支”换成“道路分支”后仍成立。
\item 每个道路连通空间连通。
\end{enumerate}
\end{proposition}

\begin{exercise}\label{exercise:appA-A-42}
证明前一个命题。
\end{exercise}

对本书中处理的大多数拓扑空间（包括所有流形），连通性与道路连通性结果是等价的。这两个概念之间的联系由下列概念提供。如果拓扑空间有一个道路连通开子集构成的基，就称它\emph{局部道路连通}。

\begin{proposition}[局部道路连通空间的性质]\label{proposition:appA-locally-path-connected}
设 $X$ 是局部道路连通拓扑空间。
\begin{enumerate}[(a)]
\item $X$ 的分支在 $X$ 中开。
\item $X$ 的道路分支等于其分支。
\item $X$ 连通，当且仅当它道路连通。
\item $X$ 的每个开子集局部道路连通。
\end{enumerate}
\end{proposition}

\begin{exercise}\label{exercise:appA-A-44}
证明前一个命题。
\end{exercise}

如果拓扑空间 $X$ 的每个开覆盖都有有限子覆盖，就称它\emph{紧}。拓扑空间的一个\emph{紧子集}是它在子空间拓扑下为紧空间的子集。例如，由 Heine--Borel 定理可知，$\mathbb R^n$ 的子集紧，当且仅当它闭且有界。

\begin{proposition}[紧空间的性质]\label{proposition:appA-compact-properties}
设 $X$ 与 $Y$ 是拓扑空间。
\begin{enumerate}[(a)]
\item 如果 $F:X\to Y$ 连续且 $X$ 紧，则 $F(X)$ 紧。
\item 如果 $X$ 紧且 $f:X\to\mathbb R$ 连续，则 $f$ 有界并在 $X$ 上达到其最大值与最小值。
\item $X$ 中有限个紧子空间之并仍紧。
\item 如果 $X$ 是 Hausdorff 的，$K,L$ 是 $X$ 的不相交紧子集，则存在不相交开子集 $U,V\subset X$ 使 $K\subset U$、$L\subset V$。
\item 紧空间的每个闭子集紧。
\item Hausdorff 空间的每个紧子集闭。
\item 度量空间的每个紧子集有界。
\item 紧空间的任意有限积紧。
\item 紧空间的任意商紧。
\end{enumerate}
\end{proposition}

\begin{exercise}\label{exercise:appA-A-46}
证明前一个命题。
\end{exercise}

对度量空间之间的映射，有几种连续性变体有用，尤其在紧空间的背景下。假设 $(M_1,d_1)$ 与 $(M_2,d_2)$ 是度量空间，$F:M_1\to M_2$ 是映射。如果对每个 $\varepsilon>0$，都存在 $\delta>0$ 使对所有 $x,y\in M_1$，$d_1(x,y)<\delta$ 蕴含 $d_2(F(x),F(y))<\varepsilon$，则称 $F$ \emph{一致连续}。如果存在常数 $C$ 使对所有 $x,y\in M_1$ 都有 $d_2(F(x),F(y))\leqslant C\,d_1(x,y)$，则称 $F$ \emph{Lipschitz 连续}。任何这样的 $C$ 称为 $F$ 的一个\emph{Lipschitz 常数}。如果每个点 $x\in M_1$ 都有一个邻域使 $F$ 在其上 Lipschitz 连续，就称 $F$ \emph{局部 Lipschitz 连续}。（为强调区别，Lipschitz 连续函数有时称为一致或整体 Lipschitz 连续。）

\begin{exercise}\label{exercise:appA-A-47}
对度量空间之间的映射，证明 Lipschitz 连续 $\Rightarrow$ 一致连续 $\Rightarrow$ 连续，且 Lipschitz 连续 $\Rightarrow$ 局部 Lipschitz 连续 $\Rightarrow$ 连续。（下面的习题 A.49 表明这些蕴含不可逆。）
\end{exercise}

\begin{proposition}\label{proposition:appA-compact-continuity}
假设 $(M_1,d_1)$ 与 $(M_2,d_2)$ 是度量空间，$F:M_1\to M_2$ 是映射。设 $K$ 是 $M_1$ 的任意紧子集。
\begin{enumerate}[(a)]
\item 如果 $F$ 连续，则 $F|_K$ 一致连续。
\item 如果 $F$ 局部 Lipschitz 连续，则 $F|_K$ Lipschitz 连续。
\end{enumerate}
\end{proposition}

\begin{proof}
先证 (a)。设 $F$ 连续，给定 $\varepsilon>0$。对每个 $x\in K$，由连续性存在正数 $\delta(x)$ 使 $d_1(x,y)<2\delta$ 蕴含 $d_2(F(x),F(y))<\varepsilon/2$。由于开球族 $\{B_{\delta(x)}(x):x\in K\}$ 覆盖 $K$，由紧性存在有限个点 $x_1,\dots,x_n\in K$ 使 $K\subset B_{\delta(x_1)}(x_1)\cup\cdots\cup B_{\delta(x_n)}(x_n)$。令 $\delta=\min\{\delta(x_1),\dots,\delta(x_n)\}$。假设 $x,y\in K$ 满足 $d_1(x,y)<\delta$。存在某个 $i$ 使 $x\in B_{\delta(x_i)}(x_i)$，于是三角不等式蕴含 $x$ 与 $y$ 都落在 $B_{2\delta(x_i)}(x_i)$ 中。由此得
\begin{align*}
d_2(F(x),F(y))\leqslant d_2(F(x),F(x_i))+d_2(F(x_i),F(y))
<\varepsilon/2+\varepsilon/2=\varepsilon.
\end{align*}

现证 (b)。设 $F$ 局部 Lipschitz 连续。由于 $F$ 连续，命题 A.45 表明 $F(K)$ 紧，因而有界。令 $D=\operatorname{diam}F(K)$。对每个 $x\in K$，存在正数 $\delta(x)$ 使 $F$ 在 $B_{2\delta(x)}(x)$ 上 Lipschitz 连续，Lipschitz 常数为 $C(x)$。由紧性，存在点 $x_1,\dots,x_n\in K$ 使 $K\subset B_{\delta(x_1)}(x_1)\cup\cdots\cup B_{\delta(x_n)}(x_n)$。令 $C=\max\{C(x_1),\dots,C(x_n)\}$，$\delta=\min\{\delta(x_1),\dots,\delta(x_n)\}$，任取 $x,y\in K$。一方面，如果 $d_1(x,y)<\delta$，则由与上一段相同的论证，$x$ 与 $y$ 落在某个使 $F$ Lipschitz 连续的球中，故 $d_2(F(x),F(y))\leqslant C\,d_1(x,y)$。另一方面，如果 $d_1(x,y)\geqslant\delta$，则 $d_2(F(x),F(y))\leqslant D\leqslant(D/\delta)d_1(x,y)$。因此 $\max\{C,D/\delta\}$ 是 $F$ 在 $K$ 上的 Lipschitz 常数。
\end{proof}

\begin{exercise}\label{exercise:appA-A-49}
设 $f,g:[0,1)\to\mathbb R$ 定义为 $f(x)=\sqrt{x}$、$g(x)=x^2$。证明 $f$ 一致连续但不局部或整体 Lipschitz 连续，而 $g$ 局部 Lipschitz 连续但不一致连续也不整体 Lipschitz 连续。
\end{exercise}

对流形、流形子集以及我们处理的大多数其他空间，紧性还有两种其他常用的等价表述。下一个命题的证明见 [LeeTM, 第 4 章]、[Mun00, 第 3 章] 和 [Sie92, 第 7 章]。

\begin{proposition}[紧性的等价表述]\label{proposition:appA-compact-equivalents}
假设 $M$ 是第二可数 Hausdorff 空间或度量空间。下列各条等价：
\begin{enumerate}[(a)]
\item $M$ 紧。
\item $M$ 的每个无限子集在 $M$ 中有极限点。
\item $M$ 中每个序列都有收敛子序列。
\end{enumerate}
\end{proposition}

\begin{exercise}\label{exercise:appA-A-51}
证明每个紧度量空间完备。
\end{exercise}

下一个引理表达紧空间最有用的性质之一。

\begin{lemma}[闭映射引理]\label{lemma:appA-closed-map-lemma}
假设 $X$ 是紧空间，$Y$ 是 Hausdorff 空间，$F:X\to Y$ 连续。
\begin{enumerate}[(a)]
\item $F$ 是闭映射。
\item 如果 $F$ 是满射，则它是商映射。
\item 如果 $F$ 是单射，则它是拓扑嵌入。
\item 如果 $F$ 是双射，则它是同胚。
\end{enumerate}
\end{lemma}

\begin{proof}
由定理 A.38，后三个断言由第一个得到，故只需证 $F$ 闭。设 $K\subset X$ 是闭子集。则命题 A.45(e) 蕴含 $K$ 紧；该命题 (a) 蕴含 $F(K)$ 紧；而 (f) 蕴含 $F(K)$ 在 $Y$ 中闭。
\end{proof}

如果 $X$ 与 $Y$ 是拓扑空间，映射 $F:X\to Y$（无论是否连续）如果对每个紧集 $K\subset Y$，原像 $F^{-1}(K)$ 紧，就称它\emph{固有}。下面是映射为固有的一些有用充分条件。

\begin{proposition}[固有性的充分条件]\label{proposition:appA-proper-sufficient}
假设 $X$ 与 $Y$ 是拓扑空间，$F:X\to Y$ 连续。
\begin{enumerate}[(a)]
\item 如果 $X$ 紧而 $Y$ Hausdorff，则 $F$ 固有。
\item 如果 $F$ 是闭映射且纤维紧，则 $F$ 固有。
\item 如果 $F$ 是闭像的拓扑嵌入，则 $F$ 固有。
\item 如果 $Y$ 是 Hausdorff 且 $F$ 有连续左逆（即连续映射 $G:Y\to X$ 使 $G\circ F=\operatorname{Id}_X$），则 $F$ 固有。
\item 如果 $F$ 固有而 $A\subset X$ 是关于 $F$ 饱和的子集，则 $F|_A:A\to F(A)$ 固有。
\end{enumerate}
\end{proposition}

\begin{exercise}\label{exercise:appA-A-54}
证明前一个命题。
\end{exercise}

\subsection*{局部紧 Hausdorff 空间}

一般而言，性质最熟悉的拓扑空间是那些拓扑由度量诱导的空间；这样的拓扑空间称为\emph{可度量化}。然而，在研究流形时，展示一个生成流形拓扑的度量往往很不方便。幸运的是，如第 1 章所示，流形属于另一类具有类似良好性质的空间——局部紧 Hausdorff 空间。本节中我们回顾这些空间的一些性质。

如果拓扑空间 $X$ 的每个点都有一个含于 $X$ 的某个紧子集中的邻域，就称它\emph{局部紧}。如果 $X$ 是 Hausdorff 的，这个性质有两个常更有用的等价表述，如下一个习题所示。

如果 $X$ 的子集在 $X$ 中的闭包紧，就称它在 $X$ 中\emph{预紧}。

\begin{exercise}\label{exercise:appA-A-55}
对 Hausdorff 空间 $X$，证明下列各条等价：
\begin{enumerate}[(a)]
\item $X$ 局部紧。
\item $X$ 的每个点都有预紧邻域。
\item $X$ 有预紧开子集构成的基。
\end{enumerate}
\end{exercise}

\begin{exercise}\label{exercise:appA-A-56}
证明局部紧 Hausdorff 空间的每个开或闭子空间本身也是局部紧 Hausdorff 空间。
\end{exercise}

下一个结果可看作闭映射引理（引理 A.52）的推广。

\begin{theorem}[固有连续映射是闭映射]\label{theorem:appA-proper-continuous-closed}
假设 $X$ 是拓扑空间，$Y$ 是局部紧 Hausdorff 空间。则每个固有连续映射 $F:X\to Y$ 是闭的。
\end{theorem}

\begin{proof}
设 $K\subset X$ 是闭子集。为证 $F(K)$ 在 $Y$ 中闭，我们证明它含有其所有极限点。设 $y$ 是 $F(K)$ 的极限点，$U$ 是 $y$ 的预紧邻域。则 $y$ 也是 $F(K)\cap\bar U$ 的极限点。由于 $F$ 固有，$F^{-1}(\bar U)$ 紧，这蕴含 $K\cap F^{-1}(\bar U)$ 紧。由于 $F$ 连续，$F(K\cap F^{-1}(\bar U))=F(K)\cap\bar U$ 紧，因而在 $Y$ 中闭。特别地，$y\in F(K)\cap\bar U\subset F(K)$，故 $F(K)$ 闭。
\end{proof}

下面是局部紧 Hausdorff 空间的一个重要性质，完备度量空间也具有该性质。证明见 [LeeTM, 第 4 章]。

\begin{theorem}[Baire 范畴定理]\label{theorem:appA-baire}
在局部紧 Hausdorff 空间或完备度量空间中，无处稠密集的可数并内部为空。
\end{theorem}

\begin{corollary}\label{corollary:appA-baire}
在局部紧 Hausdorff 空间或完备度量空间中，每个非空可数闭子集至少含一个孤立点。
\end{corollary}

\begin{proof}
设 $X$ 是这样的空间。令 $A\subset X$ 是非空可数闭子集，并假设 $A$ 没有孤立点。$A$ 在 $X$ 中闭这一事实意味着 $A$ 本身要么是局部紧 Hausdorff 空间，要么是完备度量空间。对每个 $a\in A$，单点集 $\{a\}$ 在 $A$ 中无处稠密：它在 $A$ 中闭，因为 $A$ 是 Hausdorff 的；它不含任何非空开子集，因为 $A$ 没有孤立点。由于 $A$ 是单点集的可数并，Baire 范畴定理蕴含 $A$ 在 $A$ 中内部为空，矛盾。
\end{proof}

如果在局部紧 Hausdorff 空间上再加上第二可数性这一假设，我们还能证得更多。拓扑空间 $X$ 的紧子集序列 $(K_i)_{i=1}^\infty$ 如果满足 $X=\bigcup_iK_i$ 且对每个 $i$ 都有 $K_i\subset\operatorname{Int}K_{i+1}$，就称为 $X$ 的一个\emph{紧集穷尽}。

\begin{proposition}\label{proposition:appA-exhaustion}
第二可数、局部紧 Hausdorff 空间容许一个紧集穷尽。
\end{proposition}

\begin{proof}
设 $X$ 是这样的空间。由于 $X$ 是局部紧 Hausdorff 空间，它有预紧开子集构成的基；由于它是第二可数的，它被可数个这样的集合覆盖。设 $(U_i)_{i=1}^\infty$ 是这样一个可数覆盖。从 $K_1=\bar U_1$ 开始，归纳假设已构造紧集 $K_1,\dots,K_k$ 满足对每个 $j$ 都有 $U_j\subset K_j$，且对 $j\geqslant2$ 有 $K_{j-1}\subset\operatorname{Int}K_j$。由于 $K_k$ 紧，存在某个 $m_k$ 使 $K_k\subset U_1\cup\cdots\cup U_{m_k}$。若令 $K_{k+1}=\bar U_1\cup\cdots\cup\bar U_{m_k}$，则 $K_{k+1}$ 是一个紧集，其内部含 $K_k$。此外，必要时增大 $m_k$，可设 $m_k\geqslant k+1$，从而 $U_{k+1}\subset K_{k+1}$。由归纳即得所需穷尽。
\end{proof}
\end{document}
